\begin{table}[!h]
\centering
\caption{ERPT según el signo del shock cambiario}
\centering
\begin{threeparttable}
\begin{tabular}[t]{cccccc}
\toprule
$h$ & ERPT$^{+}$ & ERPT$^{-}$ & $p$-valor & $F^{+}$ & $F^{-}$\\
\midrule
0 & -0.001 [-0.013, 0.010] & 0.004 [-0.008, 0.016] & 0.632 & --- & ---\\
3 & 0.004 [-0.039, 0.047] & 0.004 [-0.066, 0.074] & 0.952 & 54.0 & 13.4\\
6 & 0.046 [-0.056, 0.148] & 0.014 [-0.112, 0.140] & 0.638 & 51.2 & 6.0\\
12 & 0.167 [-0.190, 0.524] & 0.013 [-0.109, 0.134] & 0.522 & 5.0 & 6.8\\
18 & 0.166 [-0.198, 0.529] & 0.061 [-0.052, 0.174] & 0.891 & 2.7 & 5.1\\
\addlinespace
24 & 0.223 [-0.297, 0.742] & 0.090 [-0.009, 0.188] & 0.731 & 0.9 & 4.9\\
\bottomrule
\end{tabular}
\begin{tablenotes}
\item \textit{Note: } 
\item Intervalos al 95% con errores de Newey-West y $h+1$ rezagos. El $p$-valor corresponde al test de Wald de igualdad entre los coeficientes de respuesta del IPC ante shocks positivos y negativos. $F^{+}$ y $F^{-}$ son los estadísticos de primera etapa de cada instrumento.
\end{tablenotes}
\end{threeparttable}
\end{table}